\subsection{Retaining the alternative changes interval coverage}

The two Gaussian mechanisms in Section~\ref{exec20:expert-identification}
make it possible to examine this distinction with a known observation law.
The forward mechanism has intervention effect $3/5$ and the reverse mechanism
has effect zero, while both generate the same standardized joint distribution.
Their ordering is deliberately unresolved; this construction does not reverse
an action and an outcome whose temporal order is known.

A retained project experiment draws 2,000 independent datasets of 500
observations from that common Gaussian law. In each dataset, regression of
the outcome on an intercept and the candidate cause gives an ordinary 95\%
Student interval for the slope. Gaussian independent observations and the
correct linear conditional mean justify that slope interval. It is also an
intervention interval in the forward mechanism. Under the reverse mechanism,
the same calculation estimates a conditional association rather than the
intervention effect.

\begin{table}[htbp]
\centering\small
\begin{tabular}{lrr}
\toprule
Returned set & Forward-world coverage & Reverse-world coverage\\
\midrule
Slope interval & 95.7\% & 0.0\%\\
Slope interval together with zero & 95.7\% & 100.0\%\\
\bottomrule
\end{tabular}
\caption{One observation law and two causal interpretations. Each percentage
uses all 2,000 retained replications, each with 500 observations. The second
set retains both declared mechanisms. This is a generated example, not a
reproduction of neural-flow training or measured banking performance.}
\label{rel69:graph-coverage}
\end{table}

The slope interval covers the forward effect in 1,914 of the 2,000 datasets
and covers the reverse effect in none. Its mean width is 0.1411. The Monte
Carlo standard error of the 95.7\% coverage estimate is 0.45 percentage points;
the observed proportion does not establish improvement over nominal 95\%
coverage. First, middle and last interval calculations were independently
checked by QR decomposition in the original study. Reaggregating all retained
intervals reproduces the counts and width. This checks the stated experiment;
it supplies no calibration result for a fitted neural ensemble.

Adding zero to each slope interval produces a confidence set for the fixed
two-mechanism class. Its forward-world coverage is unchanged, and its
reverse-world coverage is one by construction. The set need not be a single
interval. Taking its convex hull would preserve coverage but obscure the
gap between the two retained causal possibilities.

This is a concrete instance of the graph-retention decomposition in
equation~\eqref{eq:la03-coverage-decomposition}. The true member is retained
with certainty because the example fixes the class to contain both mechanisms.
The forward mechanism contributes its justified regression interval; the
reverse contributes the exact zero effect. For a learned graph class, the
probability of discarding the true member must instead be controlled, and
effect inference must remain valid under the actual selection procedure.
Simultaneous inference over a fixed candidate family or an independent
estimation sample can address particular selection problems. Neither repairs
a false model or supplies missing identification. An unidentified retained
mechanism must contribute justified bounds or an uninformative set; removing
it because it widens the answer introduces an additional error.

The monetary consequence is already visible in the stated USD~100 valuation
and USD~40 intervention cost: the compatible net values are USD~20 and
minus USD~40. More precise estimation of their shared observational law does
not choose between them. A verified information clock, suitable assignment
evidence or a defensible structural restriction can change that conclusion.
The experiment therefore supports an ambiguity-preserving comparison while
leaving general graph-class coverage and bank validation open.

